\documentclass[12pt]{article}
\usepackage[utf8]{inputenc}
\usepackage[T1]{fontenc}
\usepackage{hyperref}
\title{Yabong Language Resources for K-12 Curriculum Development}
\author{Compiled by Claude Sonnet 5.5 (claude-sonnet-5-5) for Andr\'as Kornai}
\date{}
\begin{document}
\maketitle

\section{Introduction}
Yabong (ISO 639-3: ybo) is spoken in Madang Province, Papua New Guinea. Wikipedia cites about 1,500 speakers (2000 census). Other sources give very different counts: Joshua Project gives 2,500 and Global Recordings Network gives 370. Joshua Project names Baded and Mase as dialects. This search found little material for school use, and no primer, textbook or school programme.

\section{Primary Educational Resources}
The search found no primers, textbooks or school programmes for Yabong.

\section{Linguistic and Reference Materials}
\begin{itemize}
\item Wikipedia overview of the language: \url{https://en.wikipedia.org/wiki/Yabong_language}
\item ASJP wordlists (Yabong and Yabong 2): \url{https://asjp.clld.org/languages/YABONG} and \url{https://asjp.clld.org/languages/YABONG_2}
\item PARADISEC archive item (a wordlist recording): \url{https://catalog.paradisec.org.au/collections/DD1/items/1092}
\end{itemize}
Search results also mention a 2015 SIL sociolinguistic survey covering Yabong and three related languages. I did not see a direct URL for it.

\section{Religious and Community Resources}
\begin{itemize}
\item Joshua Project people-group profile: \url{https://joshuaproject.net/index.php/people_groups/15934}
\item Joshua Project language page: \url{https://joshuaproject.net/languages/ybo}
\end{itemize}
Joshua Project reports that no Bible translation is complete.

\section{Digital and Audio Resources}
\begin{itemize}
\item Global Recordings Network language page: \url{https://globalrecordings.net/en/language/ybo}
\end{itemize}
I did not check which recordings are available on that page.

\section{Curriculum Development Context}
Very little published material exists for Yabong. A curriculum would have to begin with community-led documentation: an orthography, a collection of oral texts and a primer. These could build on the existing wordlists and the survey data. Tok Pisin and English would serve as bridge languages.

\end{document}